\documentclass[11pt,a4paper]{article}
\usepackage[T1]{fontenc}
\usepackage[utf8]{inputenc}
\usepackage[polish]{babel}
\usepackage[margin=2.3cm]{geometry}
\usepackage{lmodern,microtype}
\usepackage{array,longtable,booktabs,tabularx}
\usepackage{amsmath,amssymb}
\usepackage{graphicx,xcolor,tikz,listings}
\usetikzlibrary{arrows.meta,positioning,fit,backgrounds,shapes.geometric,shapes.multipart}
\usepackage[hidelinks,unicode]{hyperref}
\usepackage{xurl}
\usepackage{bookmark}
\usepackage{enumitem}

\definecolor{accent}{HTML}{245777}
\definecolor{pale}{HTML}{EDF3F7}
\definecolor{codebg}{HTML}{F5F5F5}
\hypersetup{
  pdftitle={TSDB -- architektura systemu},
  pdfauthor={Hubert Błonowski, Jan Zaborski}
}
\setlength{\emergencystretch}{2em}
\setlength{\parskip}{0.35em}
\setlength{\parindent}{0pt}
\setcounter{tocdepth}{1}
\setlist{itemsep=0.15em,topsep=0.3em}
\renewcommand{\arraystretch}{1.18}
\setlength{\tabcolsep}{2.8pt}
\newcolumntype{P}[1]{>{\raggedright\arraybackslash}p{#1}}
\newcolumntype{Y}{>{\raggedright\arraybackslash}X}
\newcommand{\code}[1]{\expandafter\nolinkurl\expandafter{\detokenize{#1}}}
\newcommand{\tabhead}[1]{\toprule #1\\ \midrule\endfirsthead
  \toprule #1\\ \midrule\endhead \bottomrule\endfoot}
\lstset{
  basicstyle=\ttfamily\footnotesize,
  backgroundcolor=\color{codebg},
  frame=single,rulecolor=\color{black!15},
  breaklines=true,columns=fullflexible,keepspaces=true,
  showstringspaces=false,tabsize=2,
  aboveskip=0.8em,belowskip=0.8em
}
\tikzset{
  box/.style={draw=accent,rounded corners=2pt,fill=pale,
    align=center,font=\small,minimum height=11mm,text width=36mm},
  external/.style={box,fill=white,draw=black!60},
  arrow/.style={-{Latex[length=2mm]},thick,draw=accent},
  note/.style={font=\scriptsize,align=center,fill=white,inner sep=2pt},
  state/.style={box,text width=28mm,minimum height=10mm},
  umlclass/.style={draw=accent,fill=pale,rectangle split,
    rectangle split parts=3,rectangle split part align={left,left,left},
    align=left,font=\scriptsize,
    inner sep=5pt,text width=43mm},
  assoc/.style={thick,draw=accent},
  compose/.style={{Diamond[fill=accent,length=2mm]}-,thick,draw=accent},
  usecase/.style={ellipse,draw=accent,fill=pale,align=center,
    font=\small,text width=36mm,minimum height=13mm}
}


\begin{document}
\raggedbottom
\hypersetup{pageanchor=false}
\begin{titlepage}
  \centering
  {\large Politechnika Warszawska\\
    Wydział Matematyki i Nauk Informacyjnych\par}
  \vspace{2.1cm}
  {\Large Projekt zespołowy -- laboratorium 2\par}
  \vspace{1.2cm}
  {\huge\bfseries Implementacja bazy danych\\
    szeregów czasowych\\w języku Rust\par}
  \vspace{1.1cm}
  {\LARGE Architektura systemu\par}
  \vspace{1.4cm}
  \begin{tabular}{rl}
    Autorzy:          & Hubert Błonowski       \\ & Jan Zaborski\\[.5em]
    Promotor:         & mgr inż. Paweł Sobótka \\[.5em]
    Wersja dokumentu: & 1.0                    \\
    Data:             & 3 listopada 2026 r.
  \end{tabular}
  \vfill
  \url{https://github.com/janzaborski/rust-tsdb}
  \vspace{.8cm}
\end{titlepage}
\hypersetup{pageanchor=true}
\pagenumbering{roman}

\section*{Streszczenie}
Dokument przedstawia projekt jednowęzłowej bazy danych szeregów czasowych:
podział na moduły, model danych, zapis i odczyt, przechowywanie na dysku
oraz odzyskiwanie po awarii. Opisuje wykonywanie zapytań, komunikację HTTP,
integrację z Grafaną i granice dostępu do systemu.
Jest specyfikacją docelową. Parametry wydajnościowe są celami do weryfikacji.

\section*{Historia zmian}
\begin{tabularx}{\linewidth}{P{15mm}P{28mm}P{37mm}Y}
  \toprule
  Wersja & Data       & Autorzy                           & Zakres \\ \midrule
  1.0    & 03.11.2026 & H. Błonowski,\newline J. Zaborski &
  Napisanie dokumentu architektury systemu.                        \\
  \bottomrule
\end{tabularx}
\par

\begingroup
\footnotesize
\setlength{\parskip}{0pt}
\makeatletter
\renewcommand*\l@section[2]{%
  \@dottedtocline{1}{0em}{1.8em}{\bfseries #1}{#2}}
\makeatother
\tableofcontents
\endgroup
\clearpage
\pagenumbering{arabic}

\section{Cel, zakres i główne decyzje}
\label{sec:scope}
Celem jest własny silnik do przyjmowania pomiarów, przechowywania historii
i wykonywania agregacji czasowych. Użytkownik analizuje dane w Grafanie.
Zakres obejmuje HTTP, etykiety, pięć agregacji,
out-of-order w otwartym buforze oraz wdrożenie w Docker Compose.

Bloki aktywne i zamrożone umożliwiają przyjmowanie pomiarów podczas
zapisu segmentów. Poza zakresem pierwszej wersji pozostają klaster,
replikacja, kompresja, retencja, kompaktowanie, alertowanie i własny frontend.
Backup wymaga zatrzymania bazy i skopiowania katalogu danych.

\begin{figure}[!htbp]
  \centering
  \begin{tikzpicture}
    \draw[draw=accent,rounded corners] (0.6,0.6) rectangle (11,-6.3);
    \node[font=\small\bfseries,anchor=north] at (5.8,0.5)
    {TSDB z interfejsem Grafany};
    \foreach \x/\y/\id/\name in {-1/-.9/prod/Producent,-1/-4.2/anal/Analityk,
        12.4/-2.6/ops/Operator} {
        \coordinate (\id) at (\x,\y);
        \draw[thick] (\x,\y+.6) circle (.15);
        \draw[thick] (\x,\y+.45) -- (\x,\y-.12)
        (\x-.33,\y+.25) -- (\x+.33,\y+.25)
        (\x,\y-.12) -- (\x-.28,\y-.5)
        (\x,\y-.12) -- (\x+.28,\y-.5);
        \node[font=\small,below] at (\x,\y-.55) {\name};
      }
    \node[usecase] (ingest) at (3.6,-1.2) {Zapis pomiarów\\FR-01, 02, 03, 09};
    \node[usecase] (query) at (3.6,-3.4) {Zapytanie chwilowe\\lub zakresowe\\FR-04, 05, 06};
    \node[usecase] (dash) at (3.6,-5.4) {Analiza dashboardu\\FR-07};
    \node[usecase,text width=28mm] (run) at (8.6,-1.5)
    {Konfiguracja, start\\i restart bazy};
    \node[usecase,text width=28mm] (status) at (8.6,-4)
    {Kontrola zdrowia\\i gotowości};
    \draw[assoc] (prod) -- (ingest);
    \draw[assoc] (anal) -- (query);
    \draw[assoc] (anal) -- (dash);
    \draw[assoc] (ops) -- (run);
    \draw[assoc] (ops) -- (status);
  \end{tikzpicture}
  \caption{Diagram przypadków użycia UML.}
  \label{fig:usecases}
\end{figure}

\section{Architektura i odpowiedzialności modułów}
\label{sec:architecture}
Przeglądarka komunikuje się z Grafaną. Jej serwer odpytuje TSDB.
Pomiary są przyjmowane przez API HTTP. Silnik korzysta z lokalnego systemu plików.

\begin{figure}[!htbp]
  \centering
  \begin{tikzpicture}[font=\small]
    \draw[draw=accent,thick,fill=black!1] (0,0) rectangle (14,-7.7);
    \node[anchor=north west,align=left] at (.2,-.15)
    {\textit{<<device>>} Host Linux x86-64};
    \node[external,text width=40mm] (browser) at (3,-1.45)
    {Przeglądarka\\Grafana UI};
    \draw[draw=black!50,dashed] (.5,-2.8) rectangle (13.5,-5.5);
    \node[anchor=north,font=\scriptsize,align=center] at (7,-2.95)
    {\textit{<<executionEnvironment>>}\\Docker Compose, sieć tsdb-net};
    \node[box,text width=38mm] (grafana) at (3,-4.4)
    {\textit{<<container>>}\\grafana :3000};
    \node[box,text width=42mm] (db) at (10,-4.4)
    {\textit{<<container>>}\\tsdb / Rust :8080};
    % \node[external,text width=42mm] (config) at (3,-6.8)
    % {\textit{<<artifact>>}\\YAML źródła + JSON dashboardu};
    \node[external,text width=42mm] (data) at (10,-6.8)
    {Wolumen tsdb-data\\\code{/var/lib/tsdb}};
    \draw[assoc] (browser) -- node[note,right]{HTTP\\127.0.0.1:3000} (grafana);
    \draw[assoc] (grafana) -- node[note,above]{HTTP / JSON\\http://tsdb:8080} (db);
    \draw[assoc] (db.south east) -- (11.2,-5.4) --
    node[note,right]{lokalne I/O} (11.2,-6.1) -- (data.north east);
    % \draw[assoc] (config.west) -- (.7,-6.8) |-
    % node[note,pos=.6,left]{mount\\read-only} (grafana.west);
  \end{tikzpicture}
  \caption{Wdrożenie na jednym hoście.}
  \label{fig:deployment}
\end{figure}

Adapter HTTP zna protokół, parser zna składnię, a silnik zna dane
i pliki. Dzięki temu magazyn można rozwijać/testować bez HTTP,
a wykonawcę zapytań z pamięciowym interfejsem odczytu.

\begin{figure}[!htbp]
  \centering
  \begin{tikzpicture}
    \node[box,text width=95mm] (http) at (0,0)
    {Adapter HTTP: zapis, API Prometheus, stan procesu};
    \node[box] (write) at (-3,-2) {Koordynator zapisu\\podział na szeregi};
    \node[box] (query) at (3,-2) {Parser i wykonawca\\zapytań (plan)};
    \node[box,text width=108mm] (engine) at (0,-4.2)
    {Silnik: Index, indeks etykiet, Series z własnymi blokadami i blokami};
    \node[box] (wal) at (-3,-6.2) {Dziennik WAL\\zapis w tle};
    \node[box] (segment) at (3,-6.2) {Flush w tle\\segmenty i manifest};
    \draw[arrow] (http.south) -- ++(0,-0.4) -| (write.north);
    \draw[arrow] (http.south) -- ++(0,-0.4) -| (query.north);
    \draw[arrow] (write) -- (engine.north -| write);
    \draw[arrow] (query) -- (engine.north -| query);
    \draw[arrow] (engine.south -| wal) -- (wal);
    \draw[arrow] (engine.south -| segment) -- (segment);
  \end{tikzpicture}
  \caption{Zależności logiczne. }
\end{figure}

\par
\begin{tabularx}{\linewidth}{P{38mm}Y}
  \toprule
  Moduł              & Odpowiedzialność i kontrakt                                      \\ \midrule
  HTTP i walidacja   & Dekodowanie, kanonizacja, limity i błędy,
  wywołanie zapisu lub zapytania.                                                       \\
  Koordynator zapisu & Podział pracy według szeregów,
  walidacja punktów i raportowanie częściowego wyniku.                                  \\
  Index i Series     & Klucz $\leftrightarrow$ Series, dobór po etykietach,
  aktywny bufor, bloki zamrożone i segmenty dyskowe.                                    \\
  WAL                & Dziennik operacji, potwierdzanie zapisu i odtwarzanie po awarii. \\
  Magazyn segmentów  & Niezmienne pliki, katalog danych i odzyskiwanie po awarii.       \\
  Parser i wykonawca & Analiza składni, plan, odczyt i obliczenia,
  przygotowanie wyniku zapytania.                                                       \\
  \bottomrule
\end{tabularx}

\section{Model danych i współbieżność}
\label{sec:model}

\code{Index} jest katalogiem szeregów i indeksem etykiet,
a \code{Series} reprezentuje pojedynczy szereg. Są to nazwy używane w kodzie.
Klucz szeregu to \code{LabelSet}, a pomiar to \code{Sample}.
Model obejmuje bufor aktywny, niezmienne bloki oczekujące na zapis
i segmenty dyskowe. Równoległy odczyt oraz obsługa zamrożonych bloków
opisują stan docelowy

Szereg identyfikuje nazwa metryki i kanonicznie posortowane etykiety.
Wewnętrzny LabelSet zawiera także \code{__name__}.
Punkt to czas Unix w milisekundach (u64) i skończona wartość f64.
Para zajmuje 16 B. SeriesId jest lokalny dla head/WAL i może zmienić
się po restarcie. Segmenty i manifest identyfikują szeregi przez klucze.

\begin{figure}[!htbp]
  \centering
  \begin{tikzpicture}
    \node[umlclass,text width=56mm] (map) at (0,0) {
    \textbf{Index}\nodepart{second}
    inverted: HashMap<LabelSet, Arc<Series>{}>\\
    forward: HashMap<SeriesId, Arc<Series>{}>\\
    posting\_index: HashMap<String,\\
    HashMap<String, Vec<SeriesId>{}>{}>\\
    all\_ids: Vec<SeriesId>, next\_id: SeriesId
    \nodepart{third}encode(labels), get(id)\\resolve(matchers), id\_for(labels)};
    \node[umlclass,text width=45mm] (ts) at (8,0) {
      \textbf{Series}\nodepart{second}
      id: SeriesId\\labels: Arc<LabelSet>\\
      state: Mutex<SeriesState>
      \nodepart{third}new(id, labels), lock()};
    \node[umlclass,text width=50mm] (st) at (8,-3.4) {
      \textbf{SeriesState}\nodepart{second}
      samples: Vec<Sample>\\
      newest: Option<u64>\\logged\_gen: u64
      \nodepart{third}accepts(timestamp), append(sample)\\
      range(TimeRange), take\_samples()\\
      is\_logged(g), mark\_logged(g)};
    \node[umlclass,text width=43mm] (mut) at (0,-3.4) {
      \textbf{MutableBlock} (plan)\nodepart{second}
      samples: Vec<Sample> \{sorted\}
      \nodepart{third}insert(sample), slice(bounds)};
    \node[umlclass,text width=43mm] (froz) at (0,-6.3) {
    \textbf{FrozenBlock} (plan)\nodepart{second}
    samples: Vec<Sample>\\\{immutable, shared by Arc\}
    \nodepart{third}slice(bounds)};
    \node[umlclass,text width=50mm] (disk) at (8,-6.3) {
      \textbf{BlockDescriptor} (plan)\nodepart{second}
      segment, min/max, offset, count
      \nodepart{third}read(bounds)};
    \draw[assoc] (map) -- node[note,pos=.12,above]{1}
    node[note,pos=.83,above]{0..*} (ts);
    \draw[compose] (ts) -- node[note,right]{1 : 1} (st);
    \draw[compose] (st) -- node[note,above]{active: 1} (mut);
    \draw[assoc] (st.south west) -- ++(-0.8,0) |- node[note,pos=.7,above]
      {frozen: 0..1} (froz.east);
    \draw[compose] (st) -- node[note,right]{persisted: 0..*} (disk);
  \end{tikzpicture}
  \caption{Model head.}
  \label{fig:classes}
\end{figure}

Aktywny bufor utrzymuje pomiary w kolejności czasu. Opóźniony pomiar
wymaga wstawienia we właściwe miejsce. Zamrożenie przekazuje bufor
do zapisu bez pełnej kopii. Nowe pomiary trafiają do kolejnego bufora.
Widok zapytania kopiuje potrzebne dane aktywne i współdzieli niezmienne bloki.

\subsection{Współbieżność i widoczność danych}
\label{sec:concurrency}
Paczka jest jednostką transportową, a nie transakcją. Jej punkty mogą
być publikowane i odrzucane niezależnie. Błąd jednego pomiaru nie wycofuje
wcześniejszych zapisów, a odczyt może zobaczyć część przetwarzanej paczki.
Nie zapewnia się izolacji ani wspólnej migawki czasu dla wielu szeregów.

Koordynator rozdziela pracę według szeregów. Pracownik utrzymuje blokadę
tylko jednego szeregu i przetwarza małą porcję pomiarów, następnie ją zwalnia.
Nie zbiera blokad wszystkich szeregów paczki. Rejestracja nowego klucza
w katalogu jest krótka. Kodowanie, odczyt plików i oczekiwanie na
potwierdzenie dziennika nie blokują katalogu.

Zapytanie utrwala dane osobno dla każdego wybranego szeregu i zwalnia
jego blokadę przed przejściem do następnego. Poszczególne szeregi mogą
odpowiadać różnym chwilom odczytu. Szeregi utworzone w trakcie zapytania
mogą być pominięte. Jednoczesny zrzut na dysk nie może powodować
podwójnego zwracania ani pomijania już wybranych próbek.

Pozostają gwarancje poprawnego pojedynczego pomiaru, chronologicznego
odczytu i odzyskiwania potwierdzonych zapisów. Zmniejszenie zakresu blokad
ogranicza wzajemne oczekiwanie. Koszt wspólnego dziennika i obciążonych
szeregów nadal wymaga pomiaru.

\subsection{Duplikaty i opóźnione punkty}
Identyczny punkt w aktywnym buforze jest duplikatem bez zmian.
Inna wartość dla tego samego czasu odrzuca tylko ten punkt jako konflikt.
Nie modyfikuje się zamkniętych zakresów, także dla identycznego duplikatu.
Dla nowego punktu obowiązuje
\[
  t\geq H_s-W,
\]
gdzie $H_s$ to największy przyjęty czas,
a $W$ okno opóźnienia (domyślnie 5 min).
Każdy punkt porównuje się ze stanem szeregu w chwili jego przetwarzania.
W obrębie pracy dla jednego szeregu pomiary porządkuje się po czasie.
Równe czasy zachowują kolejność wejścia.

\section{Ścieżka zapisu i architektura WAL}
\label{sec:write}
\label{sec:wal-runtime}
WAL (Write-Ahead Log) pozwala odzyskać pomiary, które nie trafiły
jeszcze do segmentów danych. Po sprawdzeniu składni żądania walidacja
i rezerwacja pamięci odbywają się dla kolejnych punktów lub małych porcji
jednego szeregu. Poprawne pomiary trafiają do dziennika i są udostępniane
zapytaniom. Błędne są raportowane osobno.

Paczka HTTP może tworzyć wiele niezależnych wpisów dziennika, a pojedynczy
wpis może łączyć kilka pomiarów dla wydajności. Granica wpisu nie ustanawia
transakcji całego żądania. Awaria lub przerwanie przetwarzania może
pozostawić tylko część pomiarów paczki. Utrata połączenia nie wycofuje
już przyjętych danych.

\begin{figure}[!htbp]
  \centering
  \begin{tikzpicture}
    \foreach \x/\name in {0/HTTP,3/Zapis,6/Bufor WAL,10/Zapis w tle (pagecache),14/Synchronizacja} {
        \node[font=\small\bfseries] at (\x,.3) {\name};
        \draw[dashed,black!35] (\x,0) -- (\x,-5.4);
      }
    \draw[arrow] (0,-.5) -- node[note,above]{porcja szeregu} (3,-.5);
    \draw[arrow] (3,-1.2) -- node[note,above]{wpis do dziennika} (6,-1.2);
    \draw[arrow,dashed] (6,-1.8) -- node[note,above]{przyjęcie wpisu} (3,-1.8);
    \draw[arrow,dashed] (3,-2.45) -- node[note,above]{none: sukces} (0,-2.45);
    \node[note] at (3,-3.0) {udostępnienie danych\\zwolnienie blokad};
    \draw[arrow] (6,-3.6) -- node[note,above]{przekazanie porcji} (10,-3.6);
    \draw[arrow,dashed] (10,-4.35) -- node[note,above]{written: sukces} (0,-4.35);
    \node[note] at (14,-4.9) {okresowa synchronizacja\\bez potwierdzenia żądania};
  \end{tikzpicture}
  \caption{Przepływ zapisu małej porcji jednego szeregu. Paczka HTTP przechodzi
    ten proces wielokrotnie. Nie stanowi jednej transakcji.}
  \label{fig:write-sequence}
\end{figure}

Pracownik zapisu w tle przekazuje zebrane wpisy do pliku dziennika.
W tym czasie kolejne żądania mogą korzystać z drugiego bufora.
Osobny pracownik okresowo zleca synchronizację pliku z nośnikiem.
Rozdzielenie tych zadań pozwala łączyć małe zapisy i ogranicza czas
blokowania żądań.

\begin{tabularx}{\linewidth}{P{27mm}Y}
  \toprule
  Tryb                   & Warunek odpowiedzi po udostępnieniu danych do odczytu        \\ \midrule
  \code{none} (domyślny) & Przyjęte punkty opublikowane w pamięci. Awaria procesu
  może spowodować ich utratę.                                                           \\
  \code{written}         & Wpisy przyjętych punktów przekazane do systemu operacyjnego.
  Przyjęte dane przetrwają awarię samego procesu przy sprawnym hoście i nośniku.
  Nie ma gwarancji przetrwania utraty zasilania.                                        \\
  \bottomrule
\end{tabularx}

Okresowa synchronizacja zmniejsza ryzyko utraty danych przy awarii hosta,
ale nie jest osobnym potwierdzanym trybem żądania. Jej interwał nie
wyznacza gwarantowanego okna utraty danych: zależy ono także od obciążenia
i opóźnień urządzenia. Zapis do pamięci podręcznej systemu i trwałe
utrwalenie na nośniku są odrębnymi operacjami.

Po osiągnięciu progu rozmiaru dziennik przechodzi do kolejnego pliku.
Pojedynczy wpis nie jest dzielony między pliki. Przed checkpointem silnik
ustala zamknięty, zsynchronizowany zbiór plików. Dopiero po trwałym
opublikowaniu odpowiadających im segmentów można usunąć tę część dziennika.
Po awarii odtwarzane są pełne wpisy nowsze od checkpointu.
Błąd zapisu lub synchronizacji zatrzymuje dalsze przyjmowanie danych.
Żądania bez osiągniętego warunku potwierdzenia nie otrzymują sukcesu.

\begin{figure}[!htbp]
  \centering
  \begin{tikzpicture}
    \node[umlclass,text width=48mm] (writer) at (0,0) {
      \textbf{Interfejs WAL}\nodepart{second}porcje i tryb potwierdzenia
      \nodepart{third}przyjęcie wpisu\\oczekiwanie na zapis};
    \node[umlclass,text width=48mm] (shared) at (8,0) {
      \textbf{Stan dziennika}\nodepart{second}
      bufor wpisów\\postęp zapisu\\aktywny plik
      \nodepart{third}koordynacja pracowników};
    \node[umlclass,text width=48mm] (flush) at (0,-3.5) {
      \textbf{Pracownik zapisu}\nodepart{second}porcja danych\\bieżący plik dziennika
      \nodepart{third}zapis do systemu operacyjnego\\zamknięcie i zmiana pliku};
    \node[umlclass,text width=48mm] (sync) at (8,-3.5) {
      \textbf{Pracownik synchronizacji}\nodepart{second}ostatni postęp zapisu
      \nodepart{third}okresowe utrwalanie pliku\\na nośniku};
    \draw[assoc] (writer) -- (shared);
    \draw[arrow,dashed] (flush.north east) -- node[note,sloped,above]
    {porcja / postęp} (shared.south west);
    \draw[arrow,dashed] (shared) -- node[note,right]{aktywny plik} (sync);
    \draw[arrow,dashed] (writer) -- node[note,left]{pobudzenie} (flush);
  \end{tikzpicture}
  \caption{Elementy dziennika i ich odpowiedzialności. Zapis oraz synchronizacja
    są oddzielnymi zadaniami wykonywanymi w tle.}
  \label{fig:wal-classes}
\end{figure}

\section{Checkpoint, pliki i odzyskiwanie}
\label{sec:checkpoint}
\label{sec:storage}
Pomiary są początkowo przechowywane w pamięci i dzienniku WAL.
Co pewien czas baza zapisuje je do segmentów, czyli plików danych na dysku.
Zapis jest uruchamiany po określonym czasie (domyślnie 30 minut)
lub gdy bufor zbliża się do swojego limitu pamięci.

\subsection{Zapis danych na dysk}
\begin{enumerate}
  \item Baza wydziela dane do zapisania. Nowe pomiary trafiają do
        kolejnego bufora, dzięki czemu mogą napływać podczas zapisu.
  \item W tle powstaje segment danych. Po zakończeniu zapisu baza
        upewnia się, że plik został utrwalony na nośniku.
  \item Baza aktualizuje manifest, czyli spis gotowych segmentów
        i informację o tym, które wpisy WAL są już w nich zapisane.
        Dopiero po utrwaleniu manifestu można zwolnić odpowiedni bufor
        i usunąć niepotrzebną część WAL.
\end{enumerate}
Taki zakończony zapis tworzy checkpoint, czyli punkt, od którego
baza może rozpocząć odtwarzanie nowszych danych po awarii.
Przygotowanie danych wymaga krótkiej synchronizacji. Zapis plików
odbywa się w tle. Pojedyncza paczka HTTP może zostać zapisana częściami
w kolejnych segmentach.

\begin{figure}[!htbp]
  \centering
  \begin{tikzpicture}
    \node[anchor=east,font=\small] at (0,0) {Zapisy};
    \node[anchor=east,font=\small] at (0,-1.6) {Flush};
    \draw[arrow] (.2,0) -- (12,0);
    \draw[arrow] (.2,-1.6) -- (12,-1.6);
    \draw[fill=black!12,draw=black!40] (2,.45) rectangle (4,-2);
    \node[note] at (3,1) {przygotowanie danych\\do zapisu};
    \node[note] at (1,.3) {Stary WAL};
    \node[note] at (8,.3) {Nowe pomiary i nowy WAL};
    \draw[very thick,accent] (4.1,-1.6) -- (9.5,-1.6);
    \node[note,above] at (6.8,-1.55) {zapis segmentu i manifestu};
    \node[note] at (10.8,-2.2) {usunięcie starego WAL};
  \end{tikzpicture}
  \caption{Zapis danych na dysk odbywa się w tle, podczas przyjmowania nowych pomiarów.}
  \label{fig:checkpoint}
\end{figure}

\subsection{Odzyskiwanie po awarii}
Przy uruchomieniu baza odczytuje manifest i gotowe segmenty, a następnie
uzupełnia je pomiarami z WAL, które nie zostały jeszcze zapisane w segmentach.
Niekompletny wpis na końcu dziennika może zostać pominięty. Uszkodzenie
gotowych danych zatrzymuje uruchomienie, aby nie zwracać niepełnych wyników.
Błąd zapisu na dysk zatrzymuje dalsze przyjmowanie pomiarów.

\subsection{Buforowanie plików i kompaktowanie}
Odczyt i zapis plików korzystają z pamięci podręcznej systemu operacyjnego
(page cache).
Każdy zrzut tworzy nowy segment.

\section{Zapytania i komunikacja HTTP}
\label{sec:language}
\label{sec:http}

Istniejące endpointy to \code{POST /api/v1/write_json}
i \code{GET /api/v1/read}. Obecny zapis przyjmuje tablicę obiektów
\code{labels}/\code{samples}, z polami \code{t}/\code{v} dla pomiaru.
Poniższy kontrakt zapisu jest docelowym rozszerzeniem tego formatu.

Obsługiwany podzbiór PromQL obejmuje nazwany selektor
z dokładnym dopasowaniem etykiet, pięć funkcji czasowych i prosty skalar
(np. 1+1 do próby połączenia). Bez regexów, zagnieżdżeń i agregacji
między szeregami.
\begin{lstlisting}
temperature_celsius{service="api"}
avg_over_time(temperature_celsius{service="api"}[5m])
\end{lstlisting}
Pozostałe funkcje: \code{min_over_time}, \code{max_over_time},
\code{sum_over_time}, \code{count_over_time}.
Selektor wybiera ostatnią próbkę z $[T-D,T]$ (lookback $D=5$ min).
Funkcja działa osobno dla Series na $(T-R,T]$. Puste okno pomija Series,
również dla count. Range powtarza obliczenie dla
$T_k=start+k\cdot step\leq end$, korzystając z widoków poszczególnych szeregów.
Indeks etykiet wybiera Series, indeksy czasowe ograniczają odczyt,
a agregacje korzystają z przesuwanego okna.

\begin{figure}[!htbp]
  \centering
  \begin{tikzpicture}
    \node[box] (parse) at (0,0) {Tekst i parametry\\parser + walidacja};
    \node[box] (plan) at (5,0) {Plan\\szeregi i zakres};
    \node[box] (snap) at (10,0) {Widok szeregu\\RAM + segmenty};
    \node[box] (read) at (10,-2.2) {Odczyt zakresów\\strumień próbek};
    \node[box] (eval) at (5,-2.2) {Obliczenia\\dla siatki czasu};
    \node[box] (json) at (0,-2.2) {Wektor / macierz\\odpowiedź JSON};
    \draw[arrow] (parse) -- (plan);
    \draw[arrow] (plan) -- (snap);
    \draw[arrow] (snap) -- (read);
    \draw[arrow] (read) -- (eval);
    \draw[arrow] (eval) -- (json);
  \end{tikzpicture}
  \caption{Przepływ zapytania. Każdy etap może zakończyć się błędem.}
\end{figure}

\subsection{Architektura wykonywania zapytań}
Adapter HTTP sprawdza parametry. Parser przekształca
tekst w drzewo wyrażenia, a planista ustala szeregi, zakres odczytu
oraz siatkę czasu wyników. Indeks wybiera kandydatów według etykiet.
Metadane segmentów pozwalają pominąć pliki spoza zakresu.

Zapytanie otrzymuje widok pamięci i segmentów osobno dla każdego szeregu,
bez wspólnego punktu w czasie dla wszystkich szeregów. Odczyt może obejmować
część równocześnie przetwarzanej paczki. Blokada pojedynczego szeregu jest
zwalniana po skopiowaniu potrzebnych danych aktywnych i pobraniu uchwytów
do bloków niezmiennych, przed odczytem dyskowym i obliczeniami.
Warstwa odczytu dostarcza próbki chronologicznie, łącząc źródła bez
podwójnego zwracania punktów podczas zrzutu na dysk.

% Odczyty plików i obliczenia działają w ograniczonej puli pracowników,
% poza zadaniami obsługi sieci. Wykonawca przesuwa okno czasowe po próbkach
% i oblicza wyniki kolejnych chwil bez ponownego czytania całej historii.
% Selektor, agregacja i serializacja są oddzielnymi etapami tego samego
% wykonania; plan nie wymaga osobnego zadania dla każdej próbki.

Rozłączenie klienta sygnalizuje anulowanie wykonania. Pracownicy
sprawdzają je między porcjami odczytu i obliczeń, a następnie zwalniają
zasoby zapytania.

\subsection{Interfejs HTTP}
HTTP/1.1, UTF-8 i JSON. Zapytania także przez formularz POST.
Wartości wyników są napisami, czasy wyników sekundami Unix.
Zapis przyjmuje milisekundy i skończone f64.
\par
\begin{tabularx}{\linewidth}{P{68mm}Y}
  \toprule
  Interfejs                                                            & Główny kontrakt                         \\ \midrule
  POST \code{/api/v1/write_json}                                       & Jeden pomiar albo \code{points},
  niezależne wyniki punktów.                                                                                     \\
  GET/POST \code{/api/v1/query}                                        & query, opcjonalny time, wynik chwilowy. \\
  GET/POST \code{/api/v1/query_range}                                  & query, start, end, dodatni step,
  macierz wyników na siatce czasu.                                                                               \\
  GET \code{/api/v1/labels}, \code{/api/v1/label/<name>/values}        &
  Nazwy/wartości etykiet z całego katalogu.                                                                      \\
  GET \code{/api/v1/status/buildinfo}, \code{/healthz}, \code{/readyz} &
  Wersja TSDB, zdrowie i gotowość po recovery.                                                                   \\
  \bottomrule
\end{tabularx}
\begin{lstlisting}
POST /api/v1/write_json
{"metric":"temperature_celsius","labels":{"service":"api"},
 "timestamp":1792843200000,"value":24.5}
\end{lstlisting}
Wynik zapisu: inserted, duplicates, rejected, errors.
Zapytania stosują kopertę Prometheusa i vector/matrix/scalar.
Błędy całego żądania: 400 składnia/koperta,
422 wykonanie zapytania, 503 przeciążenie/gotowość/timeout.
Konflikty i odrzucenia pomiarów są elementami listy errors.
Timeout nie wycofuje zapisu.

\subsection{Wynik zapisu częściowego}
Błędny JSON, niepoprawna koperta, zbyt duże żądanie lub brak gotowości
stwierdzony przed rozpoczęciem przetwarzania odrzuca żądanie bez zapisu.
Po dekodowaniu błędy poszczególnych pomiarów dotyczą tylko tych punktów.
Konflikt wartości, zamknięty zakres, przekroczone okno opóźnienia albo
brak zasobów nie wycofują wcześniej przyjętych punktów.

Zakończone przetwarzanie zwraca HTTP 200 z licznikami \code{inserted},
\code{duplicates}, \code{rejected} i listą \code{errors}. Każdy błąd wskazuje
\code{point_index}, liczony od zera w oryginalnym wejściu, i \code{reason}.
Sukces obsługi żądania nie oznacza przyjęcia wszystkich pomiarów. Również
całkowite odrzucenie punktów jest rozpoznawane po licznikach.
Po pełnym przetworzeniu suma trzech liczników równa się liczbie punktów.

Tryb \code{written} potwierdza przekazanie do systemu operacyjnego
wszystkich przyjętych punktów tej odpowiedzi, także przy częściowych
odrzuceniach. Nie tworzy transakcji i nie utrzymuje blokad szeregów.
Ponowiony duplikat w otwartym buforze wymaga potwierdzenia jego
wcześniejszego zapisu, jeśli klient wybrał ten tryb.

Timeout, rozłączenie lub awaria w trakcie obsługi może pozostawić
częściowy, nieznany klientowi wynik. Błąd całego żądania nie dowodzi
wtedy braku zapisów. Ponowienie może dać duplikaty albo odrzucenia
zamkniętych zakresów. Nie gwarantuje się bezterminowej idempotencji.

\section{Technologie, wdrożenie i interfejs użytkownika}
\label{sec:technology}
\label{sec:network}
\par
\begin{tabularx}{\linewidth}{P{65mm}Y}
  \toprule
  Technologia                   & Uzasadnienie                                          \\ \midrule
  Rust jako język programowania & Ergonomia, bezpieczeństwo pamięci i kontrola zasobów. \\
  Axum, Tokio                   & Obsługa HTTP i współbieżnych żądań.                   \\
  Serde, serde\_json            & Serializacja i odczyt danych JSON.                    \\
  thiserror, tracing            & Obsługa błędów i diagnostyka.                         \\
  Grafana                       & Wizualizacja danych i zapytań.                        \\
  Docker Compose, Linux         & Powtarzalne wdrożenie na jednym hoście.               \\
  \bottomrule
\end{tabularx}

Wersje utrwala Cargo.lock, obraz Grafany jest przypięty.
Platforma docelowa: Linux x86-64 i lokalny system plików.
LOCK dopuszcza jeden proces na katalog. Usługi tsdb i grafana
korzystają z sieci tsdb-net i wolumenu tsdb-data.
Grafana odpytuje \code{http://tsdb:8080}. Przeglądarka korzysta
z portu 3000. API na porcie 8080 i Grafana są publikowane na loopback.
Baza działa bez uprawnień administratora. Profil dostępu opisano poniżej.

\subsection{Uwierzytelnianie i zakres dostępu}
Profil podstawowy obejmuje jeden zaufany host: API jest dostępne na
loopback lub w wydzielonej sieci kontenerów. Silnik nie ma własnego
systemu kont, haseł ani uprawnień do poszczególnych szeregów.
Logowanie do Grafany chroni jej interfejs. Nie zabezpiecza samodzielnie
bezpośredniego dostępu do API bazy.

Dostęp spoza zaufanego hosta wymaga osobnego profilu wdrożenia:
reverse proxy przed API zapewnia HTTPS i uwierzytelnianie, a port silnika
jest niedostępny bezpośrednio. Proponowany wariant minimalny to
Basic Auth przez TLS, obsługiwany przez gotowe proxy, np. Nginx.
Grafana potrafi przekazać dane Basic Auth
i zweryfikować certyfikat serwera.
Hasła i klucz prywatny są dostarczane jako sekrety poza repozytorium,
a certyfikat serwera musi być zaufany przez klientów.

Jest to opcjonalne rozszerzenie konfiguracji wdrożenia, do zatwierdzenia
przed udostępnieniem zdalnym. Własne zarządzanie użytkownikami,
role oraz certyfikaty klienckie pozostają poza zakresem pierwszej wersji.
Nie zakłada się, że samo dodanie hasła do Grafany zastępuje tę warstwę.

\subsection{Przykładowa konfiguracja i budżety}
Konfiguracja CLI jest wczytywana przy starcie. Compose przekazuje ją
w command. Wartości domyślne:
\par
\begin{tabularx}{\linewidth}{P{48mm}Y}
  \toprule
  Obszar          & Wartości domyślne                                   \\ \midrule
  Head i katalog  & 512 MiB active/frozen/Arc, metadata 64 MiB,
  10\,000 Series, buforowanie 30 min, out-of-order 5 min.               \\
  WAL i przyjęcie & 32 MiB na dwa bufory dziennika,
  64 MiB na ciała żądań, paczki i kodowanie,
  8 przyjętych zapisów, 4 workers.                                      \\
  Paczka          & Do 10\,000 punktów, ciało 4 MiB, wpis WAL do 8 MiB. \\
  \bottomrule
\end{tabularx}
\par
Pamięć zamrożonych bloków jest rozliczana do zwolnienia ostatniego
odczytu, także po publikacji segmentu. Zapas pojemności i zmiany
rozmiaru buforów podlegają budżetom.
Limity head nie są limitem RSS całego procesu.
Logi obejmują identyfikator żądania, wynik, czas, tryb potwierdzenia,
ilość oczekujących danych, zrzuty na dysk i raport odzyskiwania.

\subsection{Interfejs graficzny}
\label{sec:ui}
Interfejsem graficznym systemu będzie Grafana, dostępna z poziomu
przeglądarki. Użytkownik będzie mógł przeglądać pomiary na wykresach,
wybierać interesujące go szeregi i zakres czasu oraz porównywać dane
z wynikami agregacji, takimi jak średnia, minimum i maksimum.
Wykresy pozwolą obserwować trendy i zmiany wartości w czasie.
Grafana będzie pobierać dane przez API bazy, a dashboardy umożliwią
zestawienie kilku widoków danych na jednym ekranie.

% \section{Weryfikacja, ryzyka i podział pracy}
% \label{sec:verification}
% Testy obejmą częściowy zapis, liczniki błędów i kanonizację (FR-01/02),
% Written/recovery (FR-03), semantykę zapytań/funkcji (FR-04--06),
% Grafanę (FR-07), flush i out-of-order (FR-08/09).
% Wymuszone przeploty sprawdzą krótkie blokady jednego szeregu,
% widoczność części paczki, poprawność widoków szeregów,
% ponowne definicje szeregów, granice checkpointów i usuwanie dziennika.
% Awarie w Written (co najmniej 20 prób z dziennikiem klienta)
% obejmą zapis dziennika, checkpoint, manifest i usuwanie plików.

% \par
% \begin{tabularx}{\linewidth}{P{47mm}Y}
%   \toprule
%   Cel z lab1                                                            & Profil pomiaru                                         \\ \midrule
%   p95: zapis 100 ms, paczka 500 ms; odczyt 200 ms, agregacja RAM 500 ms &
%   HTTP z kolejką i Written; paczka 1000 punktów, odczyt 1000 wyników.                                                            \\
%   Co najmniej 10\,000 punktów/s                                         & 1000 Series, paczki po 1000, 4 klientów,
%   Written, godzina z flush.                                                                                                      \\
%   RSS do 1 GiB                                                          & 10\,000 punktów/s, 30 min buforowania, minimum 65 min,
%   także podczas flush i z dwoma zapytaniami.                                                                                     \\
%   Recovery do 10 s                                                      & Milion potwierdzonych punktów, osobno WAL/segmenty
%   oraz zimny/ciepły cache.                                                                                                       \\
%   \bottomrule
% \end{tabularx}

% Referencja: Linux, co najmniej 4 rdzenie, 8 GiB RAM, lokalny SSD,
% build release. Raport podaje sprzęt, konfigurację i commit;
% None mierzy się oddzielnie od trwałej przepustowości.
% Główne ryzyka to pauza podczas zamykania dziennika, pamięć odczytu i zrzutu, koszt
% out-of-order, wzrost plików bez retencji i zgodność Grafany.
% Ocenia się p95, RSS, oczekujące dane, czasy checkpointów i liczbę odrzuceń;
% cele wymagają pomiarów.

% Jan Zaborski odpowiada za WAL, segmenty, checkpoint i recovery.
% Hubert Błonowski za parser, zapytania, HTTP i Grafanę.
% Model danych, kontrakty, testy integracyjne i pomiary są wspólne.
% Wersja pełna zawiera scenariusze T-01--20,
% profile P-01--06, formaty bajtowe i komplet ustawień.

% \begingroup
% \small
% \begin{thebibliography}{10}
%   \raggedright
%   \setlength{\itemsep}{0pt}
%   \addcontentsline{toc}{section}{Źródła}
%   \bibitem{lab1} H. Błonowski, J. Zaborski,
%   \emph{Specyfikacja wymagań i zakres pracy}, lab1, wersja 1.1, 21.10.2026.
%   \bibitem{prom-api} Prometheus, \emph{HTTP API}.
%   \url{https://prometheus.io/docs/prometheus/latest/querying/api/}.
%   \bibitem{grafana-provision} Grafana Labs, \emph{Provision Grafana}.
%   \url{https://grafana.com/docs/grafana/latest/administration/provisioning/}.
%   \bibitem{linux-sync} Linux man-pages, \emph{fsync(2), fdatasync(2)}.
%   \url{https://man7.org/linux/man-pages/man2/fsync.2.html}.
%   \bibitem{linux-cache} Linux Kernel, \emph{Concepts overview: Page cache}.
%   \url{https://docs.kernel.org/admin-guide/mm/concepts.html#page-cache}.
%   \bibitem{linux-mmap} Linux man-pages, \emph{mmap(2)}.
%   \url{https://man7.org/linux/man-pages/man2/mmap.2.html}.
%   \bibitem{nginx-auth} Nginx, \emph{HTTP Basic Authentication module}.
%   \url{https://nginx.org/en/docs/http/ngx_http_auth_basic_module.html}.
%   \bibitem{nginx-tls} Nginx, \emph{Configuring HTTPS servers}.
%   \url{https://nginx.org/en/docs/http/configuring_https_servers.html}.
%   \bibitem{grafana-source} Grafana Labs, \emph{Configure the Prometheus data source}.
%   \url{https://grafana.com/docs/grafana/latest/datasources/prometheus/configure/}.
% \end{thebibliography}
% \endgroup
\end{document}
